\honest{Lotteries: A Waste of Hope}{Eliezer Yudkowsky}{2007}

The classic criticism of the lottery is that the people who play can least afford it. Some defend it, including commenters on \textsc{Overcoming Bias}, as a ``rational purchase of fantasy'': a dollar for a day of imagining yourself a millionaire.

But the fantasy has a real probability of ``nearly zero,'' that you can do nothing to bring it about, and that it is a fantasy of wealth without effort. \nb{The footnote here cites a magazine article on Carol Dweck's experiments, in which children praised for effort chose harder puzzles than children praised for intelligence. It is about praise, not about fantasies of wealth.}

So the lottery wastes hope as well as money. Without it, people ``maybe'' would dream about technical school, a business of their own or a promotion, and on the twentieth daydream ``might'' see a way to do it. I give no evidence. \nb{Oettingen and Mayer (2002) found that enjoying images of a desired future, as distinct from judging it likely, predicted less effort and worse results, in studies that included graduates looking for jobs. On that finding, the daydream the post would put in the lottery's place is not clearly better.}

I ask: ``Seriously, why can't we just say that buying lottery tickets is stupid?''

The brain cannot weaken a pleasant anticipation by a factor of 0.00000001 ``without dropping the line of reasoning entirely.'' \nb{That is close to the real odds: in 2007 a Powerball ticket had about one chance in 146 million of the jackpot. Kahneman and Tversky had written in 1979 that ``highly unlikely events are either ignored or overweighted,'' and that this favors gambling.} And people do not realize that a calculation of expected utility ``ought to override'' their instincts; they treat it as one argument among others. \nb{No calculation is given. The defenders' point is that the anticipation has value, and a calculation of expected utility would have to put a number on it.} This ``seems sufficient to explain the popularity of lotteries.'' Why, then, do so many people defend ``this classic form of self-destruction''?

Overcoming a bias takes five steps: notice it, analyze it, decide it is bad, find a workaround, and use it. Many people ``bog down in step 3,'' which ``by rights should be the easiest.'' \nb{The defenders do not agree that there is a bias to reject. That is the disagreement, not a failure to take an easy step.} ``Biases are lemons, not lemonade,'' and we should burn them.
